% =====================================================================
% preamble.tex — gói, font, màu và định dạng chung (spec v1.0)
% Biên dịch bằng XeLaTeX. Bộ màu và cách đặt tiêu đề theo báo cáo tiến độ
% giai đoạn 1: đỏ = điều khiển/trạng thái, xanh = dữ liệu ma trận.
% =====================================================================
\usepackage[a4paper,top=2.2cm,bottom=2.2cm,left=2.1cm,right=2.1cm,headheight=15pt]{geometry}

% ---- Font: TeX Gyre (đủ dấu tiếng Việt, có sẵn trong TeX Live / Overleaf)
\usepackage{fontspec}
\setmainfont{texgyreheros}[
  Extension=.otf, UprightFont=*-regular, BoldFont=*-bold,
  ItalicFont=*-italic, BoldItalicFont=*-bolditalic]
\setsansfont{texgyreheros}[
  Extension=.otf, UprightFont=*-regular, BoldFont=*-bold,
  ItalicFont=*-italic, BoldItalicFont=*-bolditalic]
\newfontfamily\headfont{texgyreadventor}[
  Extension=.otf, UprightFont=*-regular, BoldFont=*-bold,
  ItalicFont=*-italic, BoldItalicFont=*-bolditalic]
\setmonofont{DejaVuSansMono}[
  Extension=.ttf, UprightFont=*, BoldFont=*-Bold,
  ItalicFont=*-Oblique, BoldItalicFont=*-BoldOblique, Scale=0.80]
% Khối code có chú thích tiếng Việt: DejaVu Sans Mono thiếu vài dấu, dùng TeX Gyre Cursor
\newfontfamily\vnmono{texgyrecursor}[
  NFSSFamily=vnmono, Extension=.otf, UprightFont=*-regular, BoldFont=*-bold,
  ItalicFont=*-italic, BoldItalicFont=*-bolditalic, Scale=0.9]

% ---- Ngôn ngữ
\usepackage{babel}
\babelprovide[import,main]{vietnamese}
\AtBeginDocument{%
  \renewcommand{\contentsname}{Mục lục}%
  \renewcommand{\listtablename}{Danh mục bảng}%
  \renewcommand{\listfigurename}{Danh mục hình}%
  \renewcommand{\tablename}{Bảng}%
  \renewcommand{\figurename}{Hình}}

% ---- Ký hiệu toán và Unicode
\usepackage{amsmath,amssymb}
\usepackage{newunicodechar}
\newunicodechar{→}{\ensuremath{\rightarrow}}
\newunicodechar{←}{\ensuremath{\leftarrow}}
\newunicodechar{↔}{\ensuremath{\leftrightarrow}}
\newunicodechar{⇒}{\ensuremath{\Rightarrow}}
\newunicodechar{≤}{\ensuremath{\le}}
\newunicodechar{≥}{\ensuremath{\ge}}
\newunicodechar{≠}{\ensuremath{\neq}}
\newunicodechar{≈}{\ensuremath{\approx}}
\newunicodechar{∈}{\ensuremath{\in}}
\newunicodechar{∉}{\ensuremath{\notin}}
\newunicodechar{∧}{\ensuremath{\wedge}}\newunicodechar{∨}{\ensuremath{\vee}}
\newunicodechar{⌈}{\ensuremath{\lceil}}
\newunicodechar{⌉}{\ensuremath{\rceil}}
\newunicodechar{Σ}{\ensuremath{\Sigma}}
\newunicodechar{−}{\ensuremath{-}}
\newunicodechar{×}{\ensuremath{\times}}
\newunicodechar{·}{\ensuremath{\cdot}}
\newunicodechar{²}{\ensuremath{^{2}}}
\newunicodechar{³}{\ensuremath{^{3}}}

% ---- Màu (khớp slides.html của báo cáo)
\usepackage{xcolor}
\definecolor{brand}{HTML}{BC1823}
\definecolor{brandsoft}{HTML}{F6DADC}
\definecolor{brandline}{HTML}{DF9DA2}
\definecolor{cream}{HTML}{FFF5EC}
\definecolor{card}{HTML}{FFFDF9}
\definecolor{ink}{HTML}{2B2B2B}
\definecolor{muted}{HTML}{6E5F5F}
\definecolor{data}{HTML}{1D5C7A}
\definecolor{datasoft}{HTML}{E2EDF2}
\definecolor{cpu}{HTML}{4A4040}
\definecolor{okgreen}{HTML}{2E7D4F}
\color{ink}

% ---- Tiêu đề mục: chữ đỏ, vạch đỏ bên trái như tiêu đề slide
\usepackage{titlesec}
\titleformat{\section}[block]
  {\headfont\bfseries\LARGE\color{brand}}
  {\thesection}{0.6em}{}
  [{\color{brandline}\titlerule[0.8pt]}]
\titlespacing*{\section}{0pt}{14pt}{8pt}
\titleformat{\subsection}[block]
  {\headfont\bfseries\large\color{brand}}
  {\thesubsection}{0.6em}{}
\titlespacing*{\subsection}{0pt}{12pt}{4pt}
\titleformat{\subsubsection}[block]
  {\headfont\bfseries\normalsize\color{ink}}
  {\thesubsubsection}{0.6em}{}
\titlespacing*{\subsubsection}{0pt}{8pt}{2pt}
\setcounter{secnumdepth}{3}
\setcounter{tocdepth}{2}

% ---- Bảng, danh sách, code
\usepackage{array,longtable,booktabs,multirow}
\usepackage{ragged2e}
\usepackage{enumitem}
\setlist{itemsep=1pt,topsep=3pt,leftmargin=1.4em}
\usepackage{fancyvrb}
\fvset{fontsize=\small,fontfamily=vnmono}
\setlength{\LTpre}{4pt}
\setlength{\LTpost}{6pt}
\setlength{\LTcapwidth}{\linewidth}
\renewcommand{\arraystretch}{1.18}
\newcolumntype{L}[1]{>{\RaggedRight\arraybackslash}p{\dimexpr#1\linewidth-2\tabcolsep-0.6pt\relax}}
\newcolumntype{C}[1]{>{\Centering\arraybackslash}p{\dimexpr#1\linewidth-2\tabcolsep-0.6pt\relax}}

% Bảng chuẩn: \begin{spectable}{cột}{chú thích}{nhãn}{hàng tiêu đề} ... \end{spectable}
% Hàng tiêu đề lặp lại khi bảng sang trang.
\newenvironment{spectable}[4]{%
  \begingroup\small
  \begin{longtable}{@{}#1@{}}
  \caption{#2}\label{#3}\\
  \toprule #4 \\ \midrule
  \endfirsthead
  \caption[]{#2 (tiếp theo)}\\
  \toprule #4 \\ \midrule
  \endhead
  \bottomrule
  \endlastfoot}
  {\end{longtable}\endgroup}
% Bảng yêu cầu ID | Yêu cầu
\newenvironment{reqtable}[2]{%
  \begin{spectable}{L{0.12}L{0.88}}{#1}{#2}%
  {\textbf{ID} & \textbf{Yêu cầu}}}
  {\end{spectable}}

% ---- Chú thích, đánh số theo mục
\usepackage{caption}
\captionsetup{font=small,labelfont={bf,color=brand},skip=4pt}
\counterwithin{table}{section}
\counterwithin{figure}{section}

% ---- Hình (TikZ)
\usepackage{graphicx}
\usepackage{tikz}
\usetikzlibrary{arrows.meta,positioning,fit,calc,backgrounds,shapes.geometric,decorations.pathreplacing}
\tikzset{
  >={Stealth[length=2.2mm]},
  blk/.style={draw=brand,line width=0.8pt,fill=card,rounded corners=2pt,align=center,font=\small,inner sep=4pt},
  blkd/.style={draw=data,line width=0.8pt,fill=datasoft,rounded corners=2pt,align=center,font=\small,inner sep=4pt},
  blkn/.style={draw=brand,line width=1.1pt,fill=brandsoft,rounded corners=2pt,align=center,font=\small,inner sep=4pt},
  blkc/.style={draw=cpu,line width=0.8pt,fill=white,rounded corners=2pt,align=center,font=\small,inner sep=4pt},
  ctl/.style={->,draw=brand,line width=0.8pt},
  dat/.style={->,draw=data,line width=1.3pt},
  lbl/.style={font=\scriptsize,align=center,inner sep=1pt},
}
\usepackage{float}
\usepackage{placeins}
\usepackage{needspace}
\AddToHook{cmd/section/before}{\FloatBarrier\Needspace{12\baselineskip}}
\AddToHook{cmd/subsection/before}{\Needspace{10\baselineskip}}
\AddToHook{cmd/subsection/before}{\FloatBarrier}

% ---- Hộp nội dung kiểu báo cáo
\usepackage[most]{tcolorbox}
% Hộp kết luận đầu mục (một câu trả lời trước, chi tiết sau)
\newtcolorbox{keypoint}{enhanced,breakable,colback=brandsoft!55,colframe=brand,
  boxrule=0pt,leftrule=3pt,arc=0pt,left=8pt,right=6pt,top=4pt,bottom=4pt,
  fontupper=\small}
% Hộp ghi chú
\newtcolorbox{notebox}[1][Ghi chú]{enhanced,breakable,colback=datasoft!60,colframe=data,
  boxrule=0pt,leftrule=3pt,arc=0pt,left=8pt,right=6pt,top=4pt,bottom=4pt,
  fontupper=\small,title={\headfont\bfseries\footnotesize\MakeUppercase{#1}},
  coltitle=data,colbacktitle=datasoft!60,attach title to upper={\par}}
% Bốn câu hỏi của mỗi giao tiếp (như slide "Giao tiếp 1/4")
\newcommand{\qitem}[2]{{\headfont\bfseries\footnotesize\color{brand}\MakeUppercase{#1}}\par\vspace{1pt}#2\par\vspace{5pt}}
\newtcolorbox{qcard}[1]{enhanced,breakable,colback=card,colframe=brandline,boxrule=0.6pt,
  arc=2pt,left=7pt,right=7pt,top=5pt,bottom=3pt,fontupper=\small,
  title={#1},fonttitle=\headfont\bfseries\small,coltitle=brand,colbacktitle=card,
  titlerule=0.4pt,titlerule style={brandline}}
\newcommand{\interfacecard}[5]{%
  \begin{qcard}{#1}
  \qitem{Trao đổi gì?}{#2}
  \qitem{Ai khởi tạo?}{#3}
  \qitem{Khi nào được chấp nhận?}{#4}
  \qitem{Khi nào sang bước tiếp?}{#5}
  \end{qcard}}

% Nhãn trạng thái giai đoạn
\newcommand{\stagetag}[1]{\tikz[baseline=(t.base)]\node[t,fill=brand,text=white,rounded corners=1.5pt,inner sep=1.6pt,font=\headfont\bfseries\scriptsize](t){#1};}
\tikzset{t/.style={}}

% ---- Đầu trang, chân trang
\usepackage{fancyhdr}
\pagestyle{fancy}
\fancyhf{}
\fancyhead[L]{\headfont\footnotesize\color{brand}\MakeUppercase{Spec module · \ModHeader}}
\fancyhead[R]{\headfont\footnotesize\color{muted}v1.0 · spec hệ thống v1.0, gồm mục 7}
\fancyfoot[L]{\footnotesize\color{muted}Bộ tăng tốc GEMM INT8 trên Chipyard SoC}
\fancyfoot[R]{\headfont\bfseries\footnotesize\color{brand}\thepage}
\renewcommand{\headrule}{{\color{brandline}\hrule height 0.6pt}}
\renewcommand{\footrule}{{\color{brandline}\hrule height 0.6pt}}
\renewcommand{\footrulewidth}{0.6pt}
\fancypagestyle{plain}{\fancyhf{}\renewcommand{\headrulewidth}{0pt}%
  \fancyfoot[R]{\headfont\bfseries\footnotesize\color{brand}\thepage}}

\setlength{\parindent}{0pt}
\setlength{\parskip}{4pt}
\linespread{1.08}

% ---- Lệnh riêng
\newcommand{\XN}{XN}
\newcommand{\appsection}[1]{%
  \refstepcounter{section}%
  \section*{Phụ lục \thesection: #1}%
  \addcontentsline{toc}{section}{Phụ lục \thesection: #1}}

\usepackage[hidelinks]{hyperref}
\usepackage{url}
\DeclareUrlCommand\sig{\urlstyle{tt}}
\makeatletter\g@addto@macro\UrlBreaks{\do\_\do\[}\makeatother
% Bản dùng được trong chú thích và tiêu đề (đối số di chuyển)
\DeclareRobustCommand\sigc[1]{\texttt{\detokenize{#1}}}
\hypersetup{pdftitle={Spec module gemm_systolic_array v1.0},pdfauthor={Nhóm đồ án GEMM}}

\newcommand{\eqbox}[1]{\begin{tcolorbox}[enhanced,colback=card,colframe=brandline,boxrule=0.6pt,arc=2pt,left=8pt,right=8pt,top=2pt,bottom=2pt]#1\end{tcolorbox}}

% ---- Bìa dùng chung: \makecover{tên module}{tiêu đề phụ}{dòng mô tả}{các dòng bảng}
\newcommand{\makecover}[4]{%
\thispagestyle{empty}
\begin{tikzpicture}[remember picture,overlay]
  \fill[cream] (current page.north west) rectangle (current page.south east);
\end{tikzpicture}
\vspace*{3cm}
{\headfont\bfseries\small\color{brand}\MakeUppercase{Spec module · phiên bản 1.0 · full option}\par}
\vspace{0.6cm}
\begin{tikzpicture}[baseline=(b.north)]\node[fill=brand,minimum width=0.25cm,minimum height=3.6cm,inner sep=0pt](b){};\end{tikzpicture}\hspace{0.45cm}\begin{minipage}[t]{0.9\linewidth}
{\headfont\bfseries\fontsize{28}{34}\selectfont\color{brand}#1\par}
\vspace{6pt}
{\headfont\bfseries\fontsize{18}{24}\selectfont\color{brand}#2\par}
\vspace{10pt}
{\large\color{ink}#3\par}
\vspace{4pt}
{\normalsize\color{muted}Bộ tăng tốc GEMM INT8 trên Chipyard SoC\par}
\end{minipage}
\vspace{2cm}

\begingroup\small
\begin{tabular}{@{}L{0.2}L{0.78}@{}}
\toprule
#4
\bottomrule
\end{tabular}
\endgroup
\clearpage
\tableofcontents
\clearpage}


\definecolor{brand}{HTML}{000000}
\definecolor{brandline}{HTML}{888888}
\definecolor{ink}{HTML}{000000}
\color{black}
\titleformat{\section}{\bfseries\Large}{\thesection.}{0.6em}{}
\titleformat{\subsection}{\bfseries\large}{\thesubsection.}{0.6em}{}
\titleformat{\subsubsection}{\bfseries}{\thesubsubsection.}{0.6em}{}
\counterwithout{table}{section}
\counterwithout{figure}{section}
\captionsetup{labelfont=bf,justification=raggedright,singlelinecheck=false}
\fancyhead[L]{\small Đặc tả thiết kế GEMM INT8}
\fancyhead[R]{\small\ModHeader}
\fancyfoot[L]{\small Bộ tăng tốc GEMM INT8}
\fancyfoot[R]{\small\thepage}
\hypersetup{pdftitle={Đặc tả thiết kế GEMM INT8},pdfauthor={Nhóm đồ án GEMM}}

\titleformat{\section}{\color{data}\bfseries\Large}{\thesection.}{0.6em}{}
\titleformat{\subsection}{\color{data}\large}{\thesubsection.}{0.6em}{}
\renewenvironment{spectable}[4]{%
  \begingroup\small
  \begin{longtable}{#1}
  \caption{#2}\label{#3}\\
  \hline #4 \\ \hline
  \endfirsthead
  \caption[]{#2 (tiếp theo)}\\
  \hline #4 \\ \hline
  \endhead
  \endlastfoot}
  {\end{longtable}\endgroup}
\setlength{\emergencystretch}{2em}
